% Options for packages loaded elsewhere
\PassOptionsToPackage{unicode}{hyperref}
\PassOptionsToPackage{hyphens}{url}
\documentclass[
]{article}
\usepackage{xcolor}
\usepackage{amsmath,amssymb}
\setcounter{secnumdepth}{-\maxdimen} % remove section numbering
\usepackage{iftex}
\ifPDFTeX
  \usepackage[T1]{fontenc}
  \usepackage[utf8]{inputenc}
  \usepackage{textcomp} % provide euro and other symbols
\else % if luatex or xetex
  \usepackage{unicode-math} % this also loads fontspec
  \defaultfontfeatures{Scale=MatchLowercase}
  \defaultfontfeatures[\rmfamily]{Ligatures=TeX,Scale=1}
\fi
\usepackage{lmodern}
\ifPDFTeX\else
  % xetex/luatex font selection
\fi
% Use upquote if available, for straight quotes in verbatim environments
\IfFileExists{upquote.sty}{\usepackage{upquote}}{}
\IfFileExists{microtype.sty}{% use microtype if available
  \usepackage[]{microtype}
  \UseMicrotypeSet[protrusion]{basicmath} % disable protrusion for tt fonts
}{}
\makeatletter
\@ifundefined{KOMAClassName}{% if non-KOMA class
  \IfFileExists{parskip.sty}{%
    \usepackage{parskip}
  }{% else
    \setlength{\parindent}{0pt}
    \setlength{\parskip}{6pt plus 2pt minus 1pt}}
}{% if KOMA class
  \KOMAoptions{parskip=half}}
\makeatother
\setlength{\emergencystretch}{3em} % prevent overfull lines
\providecommand{\tightlist}{%
  \setlength{\itemsep}{0pt}\setlength{\parskip}{0pt}}
\usepackage{bookmark}
\IfFileExists{xurl.sty}{\usepackage{xurl}}{} % add URL line breaks if available
\urlstyle{same}
\hypersetup{
  hidelinks,
  pdfcreator={LaTeX via pandoc}}

\author{}
\date{}

\begin{document}

\section{Final Report: Design, Build, and Evaluation of a Cloud-Native
Multimodal AI
Assistant}\label{final-report-design-build-and-evaluation-of-a-cloud-native-multimodal-ai-assistant}

\subsection{Abstract}\label{abstract}

This project presents a cloud-native AI assistant deployed on GCP to
support planning, research, and task automation. The system combines a
large language model (Gemini), speech services for voice input/output,
and a tool-calling layer that can access web APIs, cloud functions, a
database, and sandboxed shell commands. The assistant is implemented as
a FastAPI backend on Cloud Run, with optional hybrid local-client usage.
Evaluation on practical test tasks shows the assistant improves speed
and convenience for common study and planning workflows compared with
manual, non-AI methods.

\subsection{1. Research Process and Tools
Explored}\label{research-process-and-tools-explored}

The research process focused on selecting a lightweight architecture
that still satisfies real cloud AI requirements.

Tools and options explored: - LLM providers: Gemini API as primary model
due to GCP alignment and strong multimodal roadmap. - Backend
frameworks: FastAPI selected for low-latency API development and simple
deployment on Cloud Run. - Speech stack: Cloud Speech-to-Text and Cloud
Text-to-Speech selected for native GCP integration. - Storage options:
Firestore chosen for schema-flexible, serverless persistence. -
Orchestration options: direct in-process tool router with optional Cloud
Function endpoint for external actions. - CI/CD and deployment: Cloud
Build + Artifact Registry + Cloud Run.

Selection rationale: - Minimal operational overhead. - Strong
cloud-native fit. - Easy extensibility for additional tools and
modalities.

\subsection{2. Technical Design}\label{technical-design}

\subsubsection{2a. High-Level Architecture
Diagram}\label{a.-high-level-architecture-diagram}

Mermaid diagram is provided in docs/architecture.mmd.

\subsubsection{2b. Cloud Services Used}\label{b.-cloud-services-used}

\begin{itemize}
\tightlist
\item
  Cloud Run: hosts assistant API service.
\item
  Cloud Functions: executes external automation actions via
  authenticated HTTP.
\item
  Firestore: stores and retrieves assistant facts/notes and persistent
  chat history (user + assistant turns).
\item
  Speech-to-Text: transcribes uploaded voice audio.
\item
  Text-to-Speech: converts generated text to audio responses.
\item
  Cloud Build and Artifact Registry: build and publish deployable
  container images.
\end{itemize}

\subsubsection{2c. Model(s) Used and Why}\label{c.-models-used-and-why}

Primary model: - Gemini via Vertex AI (gemini-2.5-flash default in
deployment): selected for low latency, strong instruction-following, and
direct GCP integration.

Optional fallback: - OpenAI model configured only as backup in
hybrid/cross-provider tests.

Why chosen: - Primary path remains fully GCP-aligned. - Fallback
improves resilience for cross-cloud interoperability experiments.

\subsubsection{2d. Data Flow and User Interaction
Flow}\label{d.-data-flow-and-user-interaction-flow}

\begin{enumerate}
\def\labelenumi{\arabic{enumi}.}
\tightlist
\item
  User sends text, voice, or image input to Cloud Run API.
\item
  Input passes security checks (length limits, sanitization, token
  checks where enabled).
\item
  Assistant planner decides whether a tool call is needed.
\item
  Tool executes (web search, Firestore lookup, Cloud Function action,
  allow-listed shell command, or allow-listed external API).
\item
  LLM generates final response based on user query plus tool output.
\item
  Response is returned as text and voice-capable output; the deployed
  homepage plays synthesized audio replies in addition to showing text.
\end{enumerate}

Homepage multimodal UX: - Single adaptive action button supports both
modalities. - Empty textbox: button acts as microphone start/end toggle.
- After capture, transcript fills textbox and button switches to send
mode. - Non-empty textbox: same button sends text. - Homepage loads
recent Firestore history for the user id so past conversation appears on
refresh.

Practical use case demonstrated: - A student asks for a weekly study
plan and requests fresh topic updates. The assistant fetches updates via
tool-calling, summarizes them, and returns a structured plan.

\subsubsection{2e. API Summary Used}\label{e.-api-summary-used}

Core assistant APIs (FastAPI endpoints): - GET /: serves the web UI. -
POST /assistant/chat: text chat with optional tool-calling. - POST
/assistant/voice/chat: end-to-end voice chat (transcribe + answer +
synthesize). - POST /assistant/voice/transcribe: speech-to-text only. -
POST /assistant/voice/synthesize: text-to-speech only. - POST
/assistant/document/analyze: uploaded document summary/extraction. - GET
/assistant/conversations: list conversation summaries. - GET
/assistant/conversations/\{session\_id\}: load one conversation thread.
- GET /assistant/history: retrieve recent stored messages.

External/cloud APIs and services integrated: - Vertex AI Gemini API
(primary LLM inference). - Google Cloud Speech-to-Text API (audio
transcription). - Google Cloud Text-to-Speech API (audio synthesis). -
Firestore API (chat and conversation persistence). - Optional Brave
Search API via web\_search tool. - Optional Cloud Function HTTP endpoint
via call\_cloud\_function tool. - Optional allow-listed external HTTP
APIs via http\_api tool.

\subsection{3. Performance Evaluation}\label{performance-evaluation}

\subsubsection{3a. Supported Task Types}\label{a.-supported-task-types}

\begin{itemize}
\tightlist
\item
  Study planning and scheduling.
\item
  Short research summarization from web results.
\item
  Basic automation via Cloud Function actions.
\item
  Voice-input conversation and speech synthesis playback in the same
  chat flow.
\end{itemize}

\subsubsection{3b. Accuracy and Utility}\label{b.-accuracy-and-utility}

Evaluation method: - Defined prompt set in evaluation/test\_cases.json.
- Automatic keyword coverage scoring in evaluation/run\_eval.py. - Human
review for relevance, structure, and actionability.

External tools available for use: - call\_cloud\_function: can be used
for automation actions through the configured function router. -
web\_search: can be used for fresh web retrieval when API key and quota
are available. - firestore\_lookup (tool): can be used to retrieve
app-stored facts/documents. - shell\_command: can be used in sandbox
mode with a strict allow-list. - http\_api: can be used for allow-listed
external API calls.

Observed results: - High utility in structured planning prompts. - Good
summarization quality when tool results are available. - Automation
responses were reliable for predefined actions.

\subsubsection{3c. Limitations and Failure
Modes}\label{c.-limitations-and-failure-modes}

\begin{itemize}
\tightlist
\item
  Tool-selection JSON may fail under ambiguous prompts.
\item
  Web-search quality depends on external API quality and quota.
\item
  Shell tool intentionally constrained; cannot perform complex commands.
\item
  Voice pipeline currently assumes standard audio format parameters.
\end{itemize}

\subsubsection{3d. Comparison to Non-AI
Workflow}\label{d.-comparison-to-non-ai-workflow}

Without assistant: - User manually searches multiple websites, organizes
notes, and executes commands separately.

With assistant: - Single interaction pipeline combines retrieval,
synthesis, and action recommendation. - Lower task-switching overhead. -
Faster first-draft planning and research summaries.

\subsection{4. Challenges and Solutions}\label{challenges-and-solutions}

Challenge 1: Securely exposing tool execution. - Solution: strict
command allow-list and syntax blocking for shell tool.

Challenge 2: Preventing API credential exposure. - Solution:
environment-based secret injection, no hard-coded keys, and
deployment-time variable controls.

Challenge 3: Keeping architecture simple but extensible. - Solution:
modular service boundaries (assistant logic, tool registry, security
module, deployment scripts).

\subsection{5. Key Learnings About Cloud AI
Systems}\label{key-learnings-about-cloud-ai-systems}

\begin{itemize}
\tightlist
\item
  Cloud-native AI assistants require both model quality and robust
  surrounding systems.
\item
  Tool-calling substantially improves practical utility beyond text
  generation alone.
\item
  Security controls are not optional; they must be designed early in
  architecture.
\item
  Serverless components reduce maintenance burden while supporting fast
  iteration.
\end{itemize}

\subsection{6. Future Improvements}\label{future-improvements}

\begin{itemize}
\tightlist
\item
  Native Gemini multimodal image caption integration in production
  endpoint.
\item
  Retrieval-augmented generation (RAG) with document indexing.
\item
  Session memory and personalization using user-scoped Firestore
  collections.
\item
  Better evaluation with factuality and latency benchmarks.
\item
  Role-based access control and Secret Manager integration for stronger
  production security.
\end{itemize}

\subsection{Security and Privacy Risks + Mitigation
Implemented}\label{security-and-privacy-risks-mitigation-implemented}

Risk 1: - API keys could leak through code commits or verbose logs.

Risk 2: - Tool-calling can expose command-injection paths if
unrestricted.

Implemented mitigation: - Input sanitization and strict command
allow-list for shell calls. - API keys moved to environment variables
and deployment configuration.

\subsection{Cross-Cloud or Hybrid Deployment
Note}\label{cross-cloud-or-hybrid-deployment-note}

This implementation supports hybrid operation: - Local client (developer
machine) sends requests to GCP-hosted assistant service. - Optional
secondary LLM provider key can be enabled for cross-provider fallback
testing.

\subsection{Reproducibility Summary}\label{reproducibility-summary}

\begin{itemize}
\tightlist
\item
  Source code structure, environment template, deployment script, and
  evaluation scripts are included.
\item
  A user with GCP project access can deploy through scripts/deploy.ps1
  and run evaluation with evaluation/run\_eval.py.
\end{itemize}

\end{document}
